The mechanism is easiest to see in three steps. Suppose first that buyers attend to energy as often for a silent listing as for a rated one. They then price silence at the expected cost of the silent pool, and the model is exactly the one estimated above. Suppose next that they attend half as often. Half of them still price silence at the expected cost of the pool, while the other half price it at whatever they implicitly assume when energy never comes up. If that default is ``nothing out of the ordinary'', the silent listing gains on every rated listing, efficient ones included. Suppose finally that buyers never attend to energy on a silent listing. It is then priced as if it carried the default cost, which can lie below the cost of an A-rated dwelling, and silence beats every rating.
